% GENERATED FILE. Edit canonical lessons, then run npm run generate:books.
% canonical-source-hash: fnv1a64:3514baea005ddff1
% canonical-lessons: ES-C689-cascada, ES-C689-arroyo, ES-C689-laguna, ES-C689-pantano, ES-C689-acantilado

\chapter{Agua en el paisaje}
\label{ch:es-agua-en-el-paisaje}
\hlchaptermodality{689}

\begin{chapteropening}
\textbf{By the end of this chapter:} \emph{I can describe water in the countryside: waterfalls, streams, lagoons, marshes and cliffs.}

Complete the last lesson of chapter 689: I can describe water in the countryside: waterfalls, streams, lagoons, marshes and cliffs.
\end{chapteropening}

\section[la cascada]{la cascada — the waterfall}

\label{lesson:ES-C689-cascada}



\begin{quote}
Before the new one: say the Spanish for the bay, then the Spanish for to float.
\end{quote}

\subsection*{You'll want to know: la cascada}
\textbf{la cascada} — \textquotedblleft{}the waterfall\textquotedblright{}.

\begin{quote}
Caminamos hasta la cascada.
\end{quote}

\begin{quote}
\emph{We walked to the waterfall.}
\end{quote}

A very big one is \emph{una catarata}.

\begin{cousinweb}
\emph{cascada} was borrowed from Italian \emph{cascata} \textquotedblleft{}a fall\textquotedblright{} (EtymDB), from \emph{cascare} \textquotedblleft{}to fall\textquotedblright{}, which goes back to Latin \emph{cadere} \textquotedblleft{}to fall\textquotedblright{}.
\end{cousinweb}

\begin{grammarlens}[title={What you've built}]
One more word for telling, explaining and describing at B1.
\end{grammarlens}

\subsection*{Your turn}
\begin{itemize}
\raggedright
  \item \emph{Say it:} \emph{la cascada}
  \item \emph{Say it:} \emph{la cascada}, once more
  \item \emph{Say it:} a sentence of your own with \emph{la cascada}
  \item \emph{Recall:} say the Spanish for the bay, then the Spanish for to float, then say \emph{la cascada} again
\end{itemize}

\begin{tcolorbox}[breakable,colback=teal!4,colframe=teal!35!black,title={Before you move on}]
What does la cascada mean? (\textquotedblleft{}The waterfall\textquotedblright{}.) Say it once more.
\end{tcolorbox}
\section[el arroyo]{el arroyo — the stream, the brook}

\label{lesson:ES-C689-arroyo}



\begin{quote}
Before the new one: say the Spanish for to float, then the Spanish for the waterfall.
\end{quote}

\subsection*{You'll want to know: el arroyo}
\textbf{el arroyo} — \textquotedblleft{}the stream, the brook\textquotedblright{}.

\begin{quote}
Cruzamos el arroyo por un puente.
\end{quote}

\begin{quote}
\emph{We crossed the stream by a bridge.}
\end{quote}

Smaller than \emph{un río}.

\begin{cousinweb}
\emph{arroyo} is usually traced to a pre-Roman word recorded in Latin as \emph{arrugia} \textquotedblleft{}a mine gallery\textquotedblright{}; the origin is not settled.
\end{cousinweb}

\begin{grammarlens}[title={What you've built}]
One more word for telling, explaining and describing at B1.
\end{grammarlens}

\subsection*{Your turn}
\begin{itemize}
\raggedright
  \item \emph{Say it:} \emph{el arroyo}
  \item \emph{Say it:} \emph{el arroyo}, once more
  \item \emph{Say it:} a sentence of your own with \emph{el arroyo}
  \item \emph{Recall:} say the Spanish for to float, then the Spanish for the waterfall, then say \emph{el arroyo} again
\end{itemize}

\begin{tcolorbox}[breakable,colback=teal!4,colframe=teal!35!black,title={Before you move on}]
What does el arroyo mean? (\textquotedblleft{}The stream, the brook\textquotedblright{}.) Say it once more.
\end{tcolorbox}
\section[la laguna]{la laguna — the lagoon, the small lake}

\label{lesson:ES-C689-laguna}



\begin{quote}
Before the new one: say the Spanish for the waterfall, then the Spanish for the stream.
\end{quote}

\subsection*{You'll want to know: la laguna}
\textbf{la laguna} — \textquotedblleft{}the lagoon, the small lake\textquotedblright{}.

\begin{quote}
Hay patos en la laguna.
\end{quote}

\begin{quote}
\emph{There are ducks on the lagoon.}
\end{quote}

\emph{Tener lagunas} — to have gaps in your memory.

\begin{cousinweb}
\emph{laguna} is inherited from Latin \emph{lacūna} \textquotedblleft{}a ditch, a pool, a gap\textquotedblright{} (EtymDB), from \emph{lacus} \textquotedblleft{}lake\textquotedblright{}.
\end{cousinweb}

\begin{grammarlens}[title={What you've built}]
One more word for telling, explaining and describing at B1.
\end{grammarlens}

\subsection*{Your turn}
\begin{itemize}
\raggedright
  \item \emph{Say it:} \emph{la laguna}
  \item \emph{Say it:} \emph{la laguna}, once more
  \item \emph{Say it:} a sentence of your own with \emph{la laguna}
  \item \emph{Recall:} say the Spanish for the waterfall, then the Spanish for the stream, then say \emph{la laguna} again
\end{itemize}

\begin{tcolorbox}[breakable,colback=teal!4,colframe=teal!35!black,title={Before you move on}]
What does la laguna mean? (\textquotedblleft{}The lagoon, the small lake\textquotedblright{}.) Say it once more.
\end{tcolorbox}
\section[el pantano]{el pantano — the marsh, the swamp; the reservoir (Spain)}

\label{lesson:ES-C689-pantano}



\begin{quote}
Before the new one: say the Spanish for the stream, then the Spanish for the lagoon.
\end{quote}

\subsection*{You'll want to know: el pantano}
\textbf{el pantano} — \textquotedblleft{}the marsh, the swamp; the reservoir (Spain)\textquotedblright{}.

\begin{quote}
Este verano el pantano tiene poca agua.
\end{quote}

\begin{quote}
\emph{This summer the reservoir has little water.}
\end{quote}

In Spain \emph{el pantano} is often a reservoir behind a dam.

\begin{cousinweb}
\emph{pantano} was borrowed from Italian \emph{pantano} \textquotedblleft{}swamp\textquotedblright{} (EtymDB); where the Italian word comes from is not settled.
\end{cousinweb}

\begin{grammarlens}[title={What you've built}]
One more word for telling, explaining and describing at B1.
\end{grammarlens}

\subsection*{Your turn}
\begin{itemize}
\raggedright
  \item \emph{Say it:} \emph{el pantano}
  \item \emph{Say it:} \emph{el pantano}, once more
  \item \emph{Say it:} a sentence of your own with \emph{el pantano}
  \item \emph{Recall:} say the Spanish for the stream, then the Spanish for the lagoon, then say \emph{el pantano} again
\end{itemize}

\begin{tcolorbox}[breakable,colback=teal!4,colframe=teal!35!black,title={Before you move on}]
What does el pantano mean? (\textquotedblleft{}The marsh, the swamp; the reservoir (Spain)\textquotedblright{}.) Say it once more.
\end{tcolorbox}
\section[el acantilado]{el acantilado — the cliff}

\label{lesson:ES-C689-acantilado}



\begin{quote}
Before the new one: say the Spanish for the lagoon, then the Spanish for the marsh.
\end{quote}

\subsection*{You'll want to know: el acantilado}
\textbf{el acantilado} — \textquotedblleft{}the cliff\textquotedblright{}.

\begin{quote}
El faro está sobre un acantilado.
\end{quote}

\begin{quote}
\emph{The lighthouse stands on a cliff.}
\end{quote}

A cliff by the sea; inland, \emph{la pared de roca}.

\begin{cousinweb}
\emph{acantilado} is built on \emph{acantilar} (EtymDB), from \emph{cantil} \textquotedblleft{}a shelf of rock\textquotedblright{}, which is built on \emph{canto} \textquotedblleft{}edge\textquotedblright{}; \emph{canto} is usually traced to Latin \emph{canthus} \textquotedblleft{}an iron tyre\textquotedblright{}.
\end{cousinweb}

\begin{grammarlens}[title={What you've built}]
One more word for telling, explaining and describing at B1.
\end{grammarlens}

\subsection*{Your turn}
\begin{itemize}
\raggedright
  \item \emph{Say it:} \emph{el acantilado}
  \item \emph{Say it:} \emph{el acantilado}, once more
  \item \emph{Say it:} a sentence of your own with \emph{el acantilado}
  \item \emph{Recall:} say the Spanish for the lagoon, then the Spanish for the marsh, then say \emph{el acantilado} again
\end{itemize}

\begin{tcolorbox}[breakable,colback=teal!4,colframe=teal!35!black,title={Before you move on}]
What does el acantilado mean? (\textquotedblleft{}The cliff\textquotedblright{}.) Say it once more.
\end{tcolorbox}
